% =====================================================================
%  fig_learning_pipeline.tex
%  Thermal Lift -- 学习流水线图（中文会议版）
%
%  独立编译：
%      xelatex fig_learning_pipeline.tex
%
%  说明：方法图的下半行。对比两种输入模式（基线 vs 混合 drizzle）
%  经过 UNet 骨干后的重建路径。关键对比：基线丢失亚像素信息，
%  混合输入保留亚像素证据并通过残差参数化锚定观测。
%  约 16cm × 6.5cm，位于 TCForge 流水线行下方。
% =====================================================================
\documentclass[border=5pt]{standalone}

\usepackage{fontspec}
\setmainfont{FandolSong-Regular}[
  Path=../fonts/,
  Extension=.otf,
  BoldFont={FandolSong-Bold},
  ItalicFont={FandolSong-Regular},
  BoldItalicFont={FandolSong-Bold},
]

\usepackage{amsmath,amssymb}
\usepackage{pifont}            % \ding{51}=checkmark, \ding{55}=cross
\usepackage{tikz}
\usetikzlibrary{positioning,arrows.meta,fit,backgrounds,calc}

% Prevent unwanted hyphenation of English terms inside nodes
\hyphenpenalty=10000
\exhyphenpenalty=10000

% ---------------------------------------------------------------- colours
\definecolor{basefill}{HTML}{F0F0F0}    % gray for baseline path
\definecolor{basedraw}{HTML}{999999}
\definecolor{propfill}{HTML}{FDF0E2}    % warm amber for proposed path
\definecolor{propdraw}{HTML}{D9772F}
\definecolor{unetfill}{HTML}{DAEAF6}    % light blue for UNet
\definecolor{unetdraw}{HTML}{4C72B0}
\definecolor{lossfill}{HTML}{FAFAFA}    % loss box
\definecolor{lossdraw}{HTML}{777777}
\definecolor{notered}{HTML}{B4474B}     % cautionary red
\definecolor{okgreen}{HTML}{2E7D32}     % affirmative green
\definecolor{divgray}{HTML}{AAAAAA}     % divider color
\definecolor{ch5color}{HTML}{C86820}    % ch5 highlight
\definecolor{lanebg1}{HTML}{F8F8F8}     % baseline lane bg
\definecolor{lanebg2}{HTML}{FFF8F0}     % proposed lane bg

% ---------------------------------------------------------------- text helpers
\newcommand{\hd}[1]{{\small\bfseries #1}}
\newcommand{\shd}[1]{{\scriptsize\bfseries #1}}   % smaller heading
\newcommand{\nt}[1]{{\scriptsize #1}}
\newcommand{\xm}{\textcolor{notered}{\ding{55}}}
\newcommand{\ok}{\textcolor{okgreen}{\ding{51}}}

% ---------------------------------------------------------------- styles
\tikzset{
  block/.style={rounded corners=2pt, draw, line width=0.6pt,
                inner sep=3pt, align=left, font=\scriptsize},
  base/.style={block, fill=basefill, draw=basedraw, text=black!70},
  prop/.style={block, fill=propfill, draw=propdraw, line width=1.0pt},
  ubox/.style={block, fill=unetfill, draw=unetdraw},
  lbox/.style={block, fill=lossfill, draw=lossdraw, densely dashed,
               rounded corners=4pt},
  resbox/.style={block, fill=propfill, draw=propdraw, line width=1.3pt},
  outbox/.style={block, fill=propfill!50, draw=propdraw},
  warnbox/.style={rounded corners=2pt, draw=notered!50, fill=notered!5,
                  line width=0.5pt, inner sep=3pt, align=left,
                  font=\scriptsize, text=notered!85},
  df/.style={-{Stealth[length=3.5pt]}, line width=0.65pt, draw=black!65},
  dfbase/.style={-{Stealth[length=3.5pt]}, line width=0.55pt,
                 draw=basedraw},
  ch5arr/.style={-{Stealth[length=4.5pt]}, line width=1.5pt,
                 draw=ch5color},
  sup/.style={-{Stealth[length=3pt]}, line width=0.5pt,
              draw=lossdraw, densely dashed},
}

\begin{document}
\begin{tikzpicture}[font=\scriptsize]

% =====================================================================
%  ROW Y COORDINATES
% =====================================================================
\def\ybase{2.5}
\def\yprop{-1.1}
\def\yloss{5.2}

% =====================================================================
%  LANE BACKGROUNDS (drawn first, behind everything)
% =====================================================================
\begin{scope}[on background layer]
  % baseline lane
  \fill[lanebg1, draw=basedraw!25, line width=0.4pt, rounded corners=3pt]
    (-1.9, 0.95) rectangle (14.1, 4.4);
  % proposed lane
  \fill[lanebg2, draw=propdraw!20, line width=0.4pt, rounded corners=3pt]
    (-1.9, -3.15) rectangle (14.1, 0.75);
\end{scope}

% =====================================================================
%  LANE LABELS (inside the lane backgrounds, top-left)
% =====================================================================
\node[font=\scriptsize\bfseries, text=basedraw, anchor=north west]
  at (-1.75, 4.3) {基线路径};
\node[font=\scriptsize\bfseries, text=propdraw, anchor=north west]
  at (-1.75, 0.6) {本文方法};

% =====================================================================
%  BASELINE PATH (top row --- grayed, represents the ablated variant)
% =====================================================================

% -- input
\node[base, text width=3.3cm] (bin) at (0, \ybase) {%
  \hd{1\texttimes{} 统计输入}\\[1pt]
  5通道：mean / median /\\
  coverage / variance / highpass\\[2pt]
  在1\texttimes{}网格构造\\
  UNet 内部 scale=2 上采样%
};

% -- UNet (baseline)
\node[base, text width=1.8cm, minimum height=1.1cm] (bunet) at (5.0, \ybase) {%
  \hd{UNet}\\[1pt]
  (scale=2)\\
  base ch=64%
};

% -- output (baseline)
\node[base, text width=1.3cm, minimum height=0.8cm] (bout) at (7.8, \ybase) {%
  \shd{输出 $\hat{y}$}\\[1pt]
  直接输出%
};

% -- failure annotation
\node[warnbox, text width=4.0cm] (bwarn) at (11.5, \ybase) {%
  \xm{} \textbf{亚像素相位在输入编码时坍塌}\\[2pt]
  1\texttimes{}网格无法编码帧间亚像素\\
  位移，网络被迫从零猜测高频细节%
};

% -- baseline flow arrows
\draw[dfbase] (bin.east) -- (bunet.west);
\draw[dfbase] (bunet.east) -- (bout.west);
\draw[dfbase, densely dotted] (bout.east) -- (bwarn.west);

% =====================================================================
%  DASHED DIVIDER between lanes
% =====================================================================
\draw[divgray, densely dashed, line width=0.7pt]
  (-1.9, 0.85) -- (14.1, 0.85);

% =====================================================================
%  PROPOSED PATH (bottom row --- highlighted, this paper's contribution)
% =====================================================================

% -- input (proposed)
\node[prop, text width=3.3cm] (pin) at (0, \yprop) {%
  \hd{Hybrid drizzle 输入}\\[1pt]
  \textbf{8通道 @ 2\texttimes{}网格}\\[2pt]
  {[ch0\textendash4]} 上采样统计通道\\
  \quad\nt{mean / median / coverage}\\
  \quad\nt{variance / highpass}\\[2pt]
  {[ch5\textendash7]} drizzle @ 2\texttimes{}:\\
  \quad\textcolor{ch5color}{\textbf{ch5: mean}} / cov / var\\[2pt]
  \ok{} \textbf{亚像素证据直接送达}%
};

% -- UNet (proposed)
\node[ubox, text width=1.8cm, minimum height=1.3cm] (punet) at (5.0, \yprop) {%
  \hd{UNet 骨干}\\[1pt]
  (scale=1)\\
  base ch=64\\[2pt]
  架构固定\\
  仅变换输入通路%
};

% -- residual head
\node[resbox, text width=3.3cm] (resid) at (8.8, \yprop) {%
  \hd{残差观测头}\\[2pt]
  $\hat{x} = \underbrace{\text{ch5}}_{\text{观测基底}} + \delta$\\[3pt]
  $\delta\!=\!0\;\Rightarrow$ 保持观测基底\\[1pt]
  惩罚：$\lambda\!\cdot\!\text{mean}|\delta|$\;(L1)\\[1pt]
  \nt{幻觉在全频域均付出代价，}\\
  \nt{无零空间可隐藏}%
};

% -- final output
\node[outbox, text width=1.5cm, minimum height=1.0cm,
      align=center] (pout) at (12.5, \yprop) {%
  \shd{2\texttimes{} 重建输出}\\[2pt]
  $\hat{x}$%
};

% -- proposed flow arrows
\draw[df] (pin.east) -- (punet.west);
\draw[df] (punet.east) --
  node[above, font=\scriptsize, text=black!55]{$\delta$}
  (resid.west);
\draw[df] (resid.east) --
  node[above, font=\scriptsize, text=black!55]{$\hat{x}$}
  (pout.west);

% -- ch5 evidence-injection arc (curves below the proposed row)
\draw[ch5arr]
  ([yshift=-2pt]pin.south east)
  to[out=-12, in=192]
  node[below, pos=0.48, font=\scriptsize, text=ch5color, yshift=-1pt]
  {ch5 drizzle mean（观测基底回加）}
  ([yshift=-2pt]resid.south west);

% =====================================================================
%  LOSS BOX (upper right corner, supervision signal)
% =====================================================================
\node[lbox, text width=3.9cm] (loss) at (8.8, \yloss) {%
  \hd{轮廓取向损失}\\
  \nt{（在合成 HR 目标上计算）}\\[2pt]
  MSE\,(0.3) + highpass\,(0.8)\\
  + SSIM\,(0.15) + 梯度向量\,(0.15)\\
  + 前向一致性边缘\,(0.05)\\[1pt]
  \nt{thin$\times$3 / gap$\times$2 加权}%
};

% -- loss supervision arrows (split from two departure points)
% To baseline output (left fork)
\draw[sup] ([xshift=-12pt]loss.south) -- ++(0,-0.25)
  -| (bout.north);
% To proposed residual head (right fork, straight down)
\draw[sup] ([xshift=8pt]loss.south) -- (resid.north);

\end{tikzpicture}
\end{document}
